\begin{apartats}

\begin{tria}{espai-tasca}
\itemtria{original}{0,75}{1,25}
Una urna conté tres boles numerades de l'1 al 3. Se'n treu una a l'atzar i, a continuació, es llança una
moneda. Descriu l'espai mostral i l'esdeveniment $A=$ «surt un nombre senar i cara».

\begin{solucio}
Cada resultat és un nombre i una cara ($C$) o una creu ($X$): $E=\{1C,\,1X,\,2C,\,2X,\,3C,\,3X\}$, amb
$3\cdot2=6$ elements.\\
$A=\{1C,\,3C\}$.
\end{solucio}

\itemtria{principi-multiplicacio}{0,75}{1,25}
Es llancen tres monedes i un dau. Quants elements té l'espai mostral? Escriu l'esdeveniment «surten
exactament dues cares i un 6», i digues quants elements té.

\begin{solucio}
Cada moneda té 2 resultats, i el dau, 6. Pel principi de multiplicació, l'espai mostral té
$2\cdot2\cdot2\cdot6=48$ elements, i no cal escriure'ls tots per saber-ho.\\
L'esdeveniment és $\{CCX6,\,CXC6,\,XCC6\}$: té 3 elements, un per a cada moneda on pot sortir la creu.
\end{solucio}
\end{tria}

\apartat[1,25]{1}
Amb els dígits 1, 2, 3, 4 i 5, quants nombres de quatre xifres diferents es poden formar? Quants són
parells? Quants comencen per 5?

\begin{solucio}
L'ordre importa i no es repeteix cap xifra: variacions sense repetició,
$V_{5,4}=5\cdot4\cdot3\cdot2=120$.\\
\textbf{Parells}: l'última xifra ha de ser 2 o 4 (2 opcions), i les altres tres són una variació de les
4 xifres que queden: $2\cdot V_{4,3}=2\cdot24=48$.\\
\textbf{Comencen per 5}: la primera xifra és fixa, i les altres tres són $V_{4,3}=24$.
\end{solucio}

\begin{nomesllarg}
\apartat{0,75}
De quantes maneres es pot triar un equip de 3 persones d'un grup de 7? I si una persona concreta hi ha
de ser?

\begin{solucio}
L'ordre no importa: combinacions, $C_{7,3}=\dbinom73=35$.\\
Si una persona hi ha de ser, només cal triar les altres dues entre les 6 que queden:
$C_{6,2}=\dbinom62=15$.
\end{solucio}
\end{nomesllarg}

\end{apartats}
